\documentclass[11pt]{article}
\usepackage{amsmath,amssymb,amsthm,booktabs}
\usepackage[margin=1.1in]{geometry}
\usepackage[numbers,sort&compress]{natbib}
\usepackage{microtype}
\usepackage{xurl}
\setlength{\emergencystretch}{1.5em}
\newtheorem{theorem}{Theorem}
\newtheorem{lemma}[theorem]{Lemma}
\newtheorem{corollary}[theorem]{Corollary}
\newtheorem{proposition}[theorem]{Proposition}
\theoremstyle{definition}
\newtheorem{definition}[theorem]{Definition}
\newcommand{\gauss}[2]{\genfrac{[}{]}{0pt}{0}{#1}{#2}_q}
\newcommand{\coeff}[1]{[q^{#1}]}
\title{\bf Sharp unimodality for a fixed-parameter slice of the Stanley--Zanello family of
Gaussian polynomial differences}
\author{Gerald\thanks{Prepared with substantial assistance from an AI system (2026);
the human author takes full responsibility for all statements and proofs. The
AI-assisted internal review recorded in the accompanying repository is not peer review.}}
\date{9 November 2026}

\begin{document}
\maketitle

\begin{abstract}
For an integer $\ell \ge 5$ let
\[
  F_\ell(q)\;=\;\gauss{\ell+8}{8}\;-\;q^{4\ell-18}\gauss{12}{6}.
\]
We prove that $F_\ell$ is unimodal for every $\ell \ge 6$, and that the bound is best possible:
$F_5$ is coefficientwise nonnegative but not unimodal. The family is a fixed-parameter slice
$(k,b)=(8,12)$ of the difference family of Stanley and Zanello, whose conjecture predicts
eventual unimodality for every fixed~$k$: Stanley and Zanello proved the cases $k\le5$, and
Koutschan--Uncu--Wong subsequently settled $k=6,7$, so $k=8$ is the smallest unresolved
parameter. The argument has three parts: an elementary
Ferrers-diagram injection gives coefficientwise nonnegativity; symmetry about the common centre
$4\ell$ reduces unimodality to nineteen fixed-offset first-difference inequalities; and those
inequalities are settled for all $\ell$ by combining an exact quasipolynomial description of
restricted partition functions with a Newton forward-difference positivity certificate, plus one
explicit finite verification for $6\le\ell\le 839$. We state precisely which steps are theorems,
which is a finite exact certificate, and what remains open.
\end{abstract}

\section{Introduction and relation to Stanley--Zanello}

Let $\gauss{n}{k}$ denote the Gaussian binomial coefficient, generating function for partitions
inside a $k\times(n-k)$ rectangle. Stanley and Zanello \cite{sz} attached to parameters
$(k,m,b)$ a difference of two Gaussian polynomials
\[
  f(k,m,b)(q)\;=\;\gauss{m}{k}-q^{s(m)}\gauss{b}{k-2},\qquad
  s(m)=\tfrac{k(m-b)}{2}+b-2k+2,
\]
admissible when $b\le(km-4k+4)/(k-2)$ and $kb\equiv km\pmod 2$, and proved nonnegativity and
unimodality completely for $k\le 5$; their Conjecture~1.1(iv) predicts the same for every $k\ge5$,
all $m\gg_k0$, and all admissible $b$ except $b=(km-4k+2)/(k-2)$ when that is an integer, and they
remark that for larger $k$ ``a new idea will most likely be necessary''. Koutschan--Uncu--Wong
\cite{kuw} settled the conjecture for $k=6$ and $k=7$ with explicit exception tables; $k\ge8$ is
open.

In this paper we settle one explicit one-parameter subfamily: with $k=8$, top parameter
$m=\ell+8$, second-summand parameter $b=12$, the shift is $s(m)=4m-50=4\ell-18$ and the difference
is exactly $F_\ell$ above. Admissibility holds precisely for $\ell\ge 5$ (the bound reads
$12\le(8\ell+36)/6$), the parity condition is automatic, and the exceptional-value condition
$b=(km-4k+2)/(k-2)$ becomes $m=51/4$, never an integer, so this slice never lies on the exceptional
locus. We do not claim the full $k=8$ case, let alone the conjecture.

Unimodality of a \emph{single} Gaussian polynomial is classical \cite{ohara,zeilberger}, and
strict unimodality is known \cite{pakpanova,zanello2015}; those results give first differences
$\ge 1$, which as we show in Section~\ref{sec:obstruction} is not enough here. Differences of
Gaussian polynomials were studied earlier by Reiner and Stanton \cite{reinerstanton}, whose
conjecture the Stanley--Zanello family generalizes. Computer-algebra resolutions of related
Stanley--Zanello families for small parameters appear in \cite{kuw}; their parameter coverage does
not include the slice treated here (see Section~\ref{sec:prior}).

\section{Definitions and notation}

All polynomials lie in $\mathbb{Z}[q]$. For $P(q)=\sum_p c_p q^p$ we write $\coeff{m}P=c_m$ and set
$c_m=0$ outside $\{0,\dots,\deg P\}$; this convention is used in every difference formula below.
$P$ is \emph{unimodal} if $c_0\le\cdots\le c_j\ge c_{j+1}\ge\cdots\ge c_D$ for some $j$, plateaus
permitted; we never claim strict unimodality. $P\ge 0$ \emph{coefficientwise} means $c_m\ge 0$ for
all $m$. For $\ell\ge 5$ put
\[
  G_\ell(q)=\gauss{\ell+8}{8},\quad a_\ell(n)=\coeff{n}G_\ell,\quad
  B(q)=\gauss{12}{6},\quad b_j=\coeff{j}B,\quad
  F_\ell=G_\ell-q^{4\ell-18}B .
\]
$\deg G_\ell=8\ell$; $\deg(q^{4\ell-18}B)=4\ell+18\le 8\ell$ exactly when $\ell\ge5$, which is the
origin of the admissible range. Both summands are palindromic about degree $4\ell$, hence so is
$F_\ell$.

\section{Main result}

\begin{theorem}\label{thm:main}
For every integer $\ell\ge 6$, $F_\ell(q)$ is unimodal.
\end{theorem}

\begin{theorem}\label{thm:nonneg}
For every integer $\ell\ge 5$, $F_\ell(q)\ge 0$ coefficientwise.
\end{theorem}

\begin{theorem}[Sharpness]\label{thm:sharp}
$F_5$ is coefficientwise nonnegative and not unimodal; explicitly
$\coeff{18..22}F_5=(15,16,15,16,15)$.
\end{theorem}

Theorem~\ref{thm:nonneg} is proved by an elementary injection (Section~\ref{sec:injection});
Theorem~\ref{thm:main} by Theorems~\ref{thm:reduction} and \ref{thm:certificate} together with
Proposition~\ref{thm:finite} below (Sections~\ref{sec:reduction}--\ref{sec:certificate});
Theorem~\ref{thm:sharp} by direct exact
computation, which is legitimate because it is a single instance.

\section{The Ferrers-diagram injection and nonnegativity}\label{sec:injection}

Identify $a_\ell(n)$ with the number of partitions $\lambda=(\lambda_1\ge\cdots\ge\lambda_8\ge0)$
of $n$ with $\lambda_1\le\ell$, and $b_j$ with partitions of $j$ inside a $6\times6$ rectangle.
Padding a partition of $4\ell-18+j$ to exactly eight parts and adding
$(1,1,1,1,0,0,0,0)$ yields a partition of $4\ell-14+j$ inside the $8\times(\ell+1)$ box: the
inequality $\lambda_4\ge\lambda_5$ is preserved strictly, the degree rises by $4=s(m+1)-s(m)$, and
the largest part rises by at most $1$, exactly the box growth. Translation by a fixed vector is
injective. Summing over the $37$ offsets $0\le j\le 36$ of the subtracted window gives the
monotonicity statement below. The exact base check is at $\ell=5$ ($m=13$): the minimum window
margin of $F_5$ is $1$, attained at degree $2$, so every window margin is at least $1\ge0$; one step
further, at $\ell=6$, the minimum window margin is $10$, attained at degree $6$. Propagating from
the base $\ell=5$ gives Theorem~\ref{thm:nonneg}. (An earlier draft of this package misstated the
base margin at $\ell=5$ as $10$; the value $10$ is the $\ell=6$ margin. Both values are recomputed
from the definitions in the accompanying code, and the verification program asserts them.)

\begin{theorem}[Monotone window injection, even $k$ \cite{sz}]\label{thm:injection}
Let $k\ge4$ be even and $(m,b)$ admissible, and let
$f(k,m,b)=\gauss{m}{k}-q^{s(m)}\gauss{b}{k-2}$ with
$s(m)=k(m-b)/2+b-2k+2$. Then for every fixed offset $j\ge0$,
\[
  \coeff{s(m)+j}\gauss{m}{k}\;\le\;\coeff{s(m+1)+j}\gauss{m+1}{k}.
\]
Consequently coefficientwise nonnegativity of $f(k,m,b)$ is monotone in $m$: once it holds at an
admissible $m_0$ it holds for all admissible $m\ge m_0$, so eventual nonnegativity for fixed
$(k,b)$ reduces to a finite exact check.
\end{theorem}

\begin{proof}[Proof of Theorem~\ref{thm:injection}]
Let $k\ge4$ be even and let $(m,b)$ be admissible: $3\le k\le m$, $3\le b\le(km-4k+4)/(k-2)$, and
$kb\equiv km\pmod2$; then $s(m)=k(m-b)/2+b-2k+2$ is an integer, and admissibility of $(m+1,b)$ is
automatic because $k$ is even and the upper bound on $b$ is nondecreasing in $m$. Represent
$\gauss{m}{k}$ as generating partitions $\lambda=(\lambda_1\ge\cdots\ge\lambda_k\ge0)$ of the box
$k\times(m-k)$, that is, with $\lambda_1\le m-k$, and degree $|\lambda|$. Define
\[
  \Phi(\lambda)=\bigl(\lambda_1+1,\dots,\lambda_{k/2}+1,\ \lambda_{k/2+1},\dots,\lambda_k\bigr).
\]
Then $\Phi(\lambda)$ is a partition: the only adjacent pair affected is $(k/2,\,k/2+1)$, where
$\lambda_{k/2}+1\ge\lambda_{k/2+1}+1>\lambda_{k/2+1}$; equivalently, the Ferrers diagram grows by
exactly one cell in each of its first $k/2$ rows, which preserves the partition shape because row
lengths are weakly decreasing. It lies in the $k\times(m+1-k)$ box because
$\lambda_1+1\le(m-k)+1=(m+1)-k$. Its degree is $|\lambda|+k/2$, and $s(m+1)-s(m)=k/2$ since $k$ is
even. Hence $\Phi$ maps partitions of $s(m)+j$ in the $k\times(m-k)$ box into partitions of
$s(m+1)+j$ in the $k\times(m+1-k)$ box; it is injective because it is translation by a fixed
vector. Comparing coefficients yields $\coeff{s(m)+j}\gauss{m}{k}\le
\coeff{s(m+1)+j}\gauss{m+1}{k}$ for every fixed offset $j\ge0$.

For nonnegativity, outside the window $\{s(m),\dots,s(m)+\deg\gauss{b}{k-2}\}$ the coefficients of
$f(k,m,b)$ equal those of $\gauss{m}{k}$ and are nonnegative; inside, the coefficient at offset $j$
is $\coeff{s(m)+j}\gauss{m}{k}-b_j$, where $b_j=\coeff{j}\gauss{b}{k-2}$ does not depend on $m$. The
inequality above shows each window margin is nondecreasing along admissible $m$, so nonnegativity at
one admissible $m_0$ --- a finite check of $\deg\gauss{b}{k-2}+1$ integers --- implies
nonnegativity for all admissible $m\ge m_0$. For the slice $(k,b)=(8,12)$ the base is $m_0=13$
($\ell=5$) with minimum window margin $1$, as computed above.
\end{proof}

The argument requires $k$ even; the odd-$k$ analogue is not claimed. The injection controls
nonnegativity only, and its technique is adjacent to the column-growth step of O'Hara's proof
\cite{ohara}; see Section~\ref{sec:prior} for the priority discussion.

\section{Adjacent-difference reduction}\label{sec:reduction}

Because $F_\ell$ is symmetric about $4\ell$, unimodality is equivalent to
$\coeff{d}F_\ell\le\coeff{d+1}F_\ell$ for $0\le d\le 4\ell-1$. Writing $j=d-(4\ell-18)$,
\[
  \coeff{d+1}F_\ell-\coeff{d}F_\ell
  \;=\;\bigl(a_\ell(d+1)-a_\ell(d)\bigr)-\bigl(b_{j+1}-b_j\bigr).
\]
For $j\le-2$ both $b$-indices are negative and the right-hand side equals
$a_\ell(d+1)-a_\ell(d)\ge0$ by unimodality of $G_\ell$, since $d+1<4\ell$. For $j\ge18$ symmetry
applies. Exactly $j=-1,0,\dots,17$ carry conditions. The endpoint $j=-1$ is not free:
$b_0-b_{-1}=1$.

\begin{theorem}[Reduction]\label{thm:reduction}
For every $\ell\ge6$, $F_\ell$ is unimodal if and only if
\[
  a_\ell(4\ell-18+j+1)-a_\ell(4\ell-18+j)\;\ge\; D b_j
  \qquad (j=-1,0,\dots,17),
\]
where $D b_j=b_{j+1}-b_j$ is the vector displayed in Table~\ref{tab:db}.
\end{theorem}

\begin{table}[h]
\centering
\begin{tabular}{crrrrrrrrrrrrrrrrrrr}
\toprule
$j$ & $-1$&$0$&$1$&$2$&$3$&$4$&$5$&$6$&$7$&$8$&$9$&$10$&$11$&$12$&$13$&$14$&$15$&$16$&$17$\\
\midrule
$Db_j$ & $1$&$0$&$1$&$1$&$2$&$2$&$4$&$2$&$5$&$4$&$6$&$4$&$7$&$3$&$6$&$3$&$4$&$0$&$3$\\
\bottomrule
\end{tabular}
\caption{The nineteen required lower bounds, computed from $\gauss{12}{6}$ in the accompanying
code rather than tabulated by hand. The interior maximum is $Db_{11}=7$; the binding case for large
$\ell$ is $j=17$, the step immediately before the centre.}
\label{tab:db}
\end{table}

\section{Quasipolynomial certification}\label{sec:certificate}

\begin{lemma}\label{lem:quasi}
Let $p(n)=\coeff{n}\prod_{i=1}^{8}(1-q^i)^{-1}$ with $p(n)=0$ for $n<0$. Then $p$ is a
quasipolynomial in $n$ of degree at most $7$ and period dividing $840=\mathrm{lcm}(1,\dots,8)$,
valid for every $n\ge0$.
\end{lemma}

\begin{proof}
The generating function $P(q)=1/D(q)$ with $D(q)=\prod_{i=1}^{8}(1-q^i)$ is a \emph{proper} rational
function: the numerator is $1$ and $\deg D=1+\cdots+8=36>0$. Its poles are the roots of unity
$\zeta$ of order dividing $840=\mathrm{lcm}(1,\dots,8)$, and the multiplicity of $\zeta$ as a pole
is $m(\zeta)=|\{1\le i\le 8: \zeta^i=1\}|\le 8$, attained only at $\zeta=1$. Because the fraction
is proper, its partial fraction decomposition over $\mathbb{C}$ has \emph{no polynomial part}: it is
the exact identity of rational functions
\[
  \frac{1}{D(q)}
  \;=\;\sum_{\zeta^{840}=1}\ \sum_{r=1}^{m(\zeta)}
        \frac{c_{\zeta,r}}{(1-q/\zeta)^{r}},
  \qquad c_{\zeta,r}\in\mathbb{C},
\]
summed over all $840$-th roots of unity. Expanding each term as a geometric-type series,
$\coeff{n}(1-q/\zeta)^{-r}=\binom{n+r-1}{r-1}\zeta^{-n}$, an identity of formal power series valid
for \emph{every} $n\ge0$. Hence for all $n\ge0$,
\[
  p(n)=\sum_{\zeta^{840}=1}\ \sum_{r=1}^{m(\zeta)}
        c_{\zeta,r}\binom{n+r-1}{r-1}\zeta^{-n},
\]
and since $\zeta^{-840}=1$, on each residue class $n\equiv\rho\pmod{840}$ the factors $\zeta^{-n}$
are constant, so $p(\rho+840t)$ is a polynomial in $t$ (equivalently in $n$) of degree at most
$\max_\zeta m(\zeta)-1=7$. Grouping conjugate summands shows the class polynomials have rational
coefficients; they take integer values because $p$ does. No ``sufficiently large $n$'' qualification
appears anywhere: the only finitely-supported contributions that ordinary asymptotic quasipolynomial
statements need are absent here precisely because the numerator is constant and the decomposition
carries no polynomial part.
\end{proof}

\begin{lemma}[Inclusion--exclusion]\label{lem:ie}
$a_\ell(n)=\sum_{S\subseteq\{1,\dots,8\}}(-1)^{|S|}\,p\bigl(n-|S|\ell-\sigma(S)\bigr)$,
$\sigma(S)=\sum_{i\in S}i$.
\end{lemma}

\begin{proof}
Expand $\prod_{i=1}^{8}(1-q^{\ell+i})$ in
$\gauss{\ell+8}{8}=\prod_{i=1}^{8}(1-q^{\ell+i})\prod_{i=1}^{8}(1-q^i)^{-1}$ and extract
coefficients.
\end{proof}

For the offsets at hand $n=4\ell-18+j$ with $j\le17$, so the argument of $p$ belonging to a subset
$S$ with $|S|=4$ is at most $-1-10=-11<0$, and for $|S|\ge5$ it is at most $-\ell-11$. Hence:

\begin{corollary}\label{cor:poly}
For $\ell\ge840$ and each offset $j$, the margin
$Q_{j,r}(t):=a_\ell(4\ell-18+j+1)-a_\ell(4\ell-18+j)-Db_j$ with
$\ell=r+840(t+1)$, $0\le r<840$, is a polynomial in $t$ of degree at most $7$, valid for all
integer $t\ge0$.
\end{corollary}

\begin{lemma}[Newton positivity]\label{lem:newton}
If a real polynomial $Q$ of degree at most $7$ satisfies $\Delta^i Q(0)\ge0$ for $i=0,\dots,7$,
where $\Delta$ is the forward difference, then $Q(t)\ge0$ for every integer $t\ge0$.
\end{lemma}

\begin{proof}
Newton's expansion $Q(t)=\sum_{i=0}^{7}\binom{t}{i}\Delta^iQ(0)$ with $\binom{t}{i}\ge0$ for
integer $t\ge0$.
\end{proof}

\begin{theorem}[Certificate]\label{thm:certificate}
For every $j\in\{-1,\dots,17\}$ and every $r\in\{0,\dots,839\}$, the eight Newton coefficients
$\Delta^i Q_{j,r}(0)$, $i=0,\dots,7$, are nonnegative. Consequently the nineteen inequalities of
Theorem~\ref{thm:reduction} hold for all $\ell\ge840$.
\end{theorem}

The theorem's hypothesis is a finite assertion about $19\cdot840=15{,}960$ integer octuples,
produced and validated by the accompanying code (Section~\ref{sec:repro}). Two features are worth
recording. First, the sixth and seventh Newton coefficients vanish identically in every class, so
the true degree is at most $5$; the degree-$7$ bound is used only to size the certificate. Second,
the smallest tail margin is $3{,}552{,}106{,}826$, so the tail is far from sharp: all the tension in
Theorem~\ref{thm:main} lives in the finite base $6\le\ell\le839$, where the minimum margin is $1$,
attained at $(\ell,j)=(6,16)$, and where the central surplus equals exactly $4$ against the required
$3$ at $\ell=6,7$.

The finite base is likewise a formal component of the proof chain, not an afterthought.

\begin{proposition}[Finite base]\label{thm:finite}
For every $6\le\ell\le839$ and every $j\in\{-1,0,\dots,17\}$, the first-difference inequality of
Theorem~\ref{thm:reduction} holds. The verification is exact integer computation over all
$19\cdot834$ instances (no floating point); the minimum margin is~$1$, attained at
$(\ell,j)=(6,16)$.
\end{proposition}

\begin{proof}
For each $\ell$ in the range, the coefficients $a_\ell(n)$ for $n\le4\ell$ are computed exactly as
follows: the power series of $\prod_{i=1}^{8}(1-q^{i})^{-1}$ is built by the coin-change recurrence
and truncated at degree $4\ell$, then the numerator factors $1-q^{\ell+i}$, $i=1,\dots,8$, are
applied by in-place descending subtraction. Since every numerator factor has nonnegative exponents,
the truncation at degree $4\ell$ is exact for all coefficients accessed, and the result is
$[\ell+8\choose 8]_q$ up to degree $4\ell$. The first differences
$a_\ell(4\ell-18+j+1)-a_\ell(4\ell-18+j)$ are then compared with $Db_j$ in exact integers.
The accompanying package performs this computation through two distinct code paths (the
inclusion--exclusion route of the generator, Lemma~\ref{lem:ie}, and the dense truncated
multiplication route of the verifier and \texttt{r1\_decisive.py}); both report the same minimum
margin and its location.
\end{proof}

\subsection{Why this is a certificate and not experimental evidence}\label{sec:why}
Corollary~\ref{cor:poly} and Lemma~\ref{lem:newton} together turn the computation into a proof of a
statement about \emph{all} $\ell\ge840$: the computation only verifies the hypotheses of a theorem
about polynomials, over an explicitly enumerated finite set. No interpolation, no floating point, and
no asymptotic fitting is involved; the values are integers obtained from the definitions, and the
verification package recomputes the complete Newton vectors of every obligation through a second
independent code path (dense truncated Gaussian multiplication), so a value-tampered certificate
cannot pass.

\subsection{The obstruction to an injection-only proof}\label{sec:obstruction}
The injection of Section~\ref{sec:injection} can be sharpened to an exact identity. Adding $1$ to
the largest part maps partitions of $n$ with $\lambda_1\le\ell-1$ bijectively onto partitions of
$n+1$ with $\mu_1\ge\mu_2+1$, whence
\[
  a_\ell(n+1)-a_\ell(n)
  \;=\;\#\{\mu\vdash n+1 \text{ in the } 8\times\ell \text{ box}: \mu_1=\mu_2\}
   \;-\;\#\{\lambda\vdash n \text{ in the box}: \lambda_1=\ell\}.
\]
Both terms are, at fixed offset, degree-$6$ quasi-polynomials in $\ell$, and their degree-$6$
leading terms cancel exactly while the degree-$5$ term survives: the difference is of degree
exactly $5$ on each residue class modulo $840$---precisely what the certificate records (the fifth
Newton coefficient is nonzero in all $15{,}960$ classes and the sixth and seventh vanish
identically). The mechanism is the palindromic symmetry $a_\ell(m)=a_\ell(8\ell-m)$: the window
sits at the centre of the coefficient array, so the continuous slice-volume expansion in powers
of $1/\ell$ is even about the centre and its first difference contributes no degree-$7$ or
degree-$6$ term, with leading smooth term proportional to $(c+1)^2-c^2=2c+1$ at degree $\ell^5$
(the lower-degree periodic contributions are unrestricted but cannot raise the degree above $5$);
exact rational interpolation on sampled residue classes confirms that the
degree-$7$ and degree-$6$ terms of $a_\ell(4\ell+c)$ agree with those of $a_\ell(4\ell+c+1)$
(hence cancel in the difference) and that the surviving degree-$5$ coefficient is proportional to
$2c+1$, so it is nonzero for every integer offset $c$. A residual of degree $5$ must be compared
against absolute constants ($3$ and $7$). An injection yields $\ge0$; here one needs $\ge3$, and equality is nearly attained
(true value $4$ at $\ell=6$). Box symmetries (conjugation, complement) supply no surplus set of the
required size. This is the precise sense in which the elementary mechanism does not extend to
Theorem~\ref{thm:main}, and why a quantitative input of a different kind is needed.

\section{Finite certificate and reproducibility}\label{sec:repro}

The companion directory \texttt{verification/} contains three standard-library Python programs and a
recorded expectation file:
\texttt{generate\_certificate.py} builds the $15{,}960$-obligation certificate from the definitions
(offset constants derived from $\gauss{12}{6}$ inside the program, never read from a table);
\texttt{verify\_certificate.py} re-derives those constants, checks structural completeness (no
missing or duplicated obligations), checks nonnegativity of all Newton coefficients, checks the
degree drop, and recomputes the full eight-coefficient Newton vector of \emph{every} one of the
$15{,}960$ obligations through an independent dense code path (plus a seeded sample through a
third route), failing non-zero on a missing, truncated, or value-tampered certificate;
\texttt{independent\_check.py} verifies the theorem statements themselves by constructing $F_\ell$
completely for $\ell=5,6,7$, sweeping $6\le\ell\le80$, testing large parameters, and confirming the
equivalence of Theorem~\ref{thm:reduction} against direct unimodality for $5\le\ell\le60$. Reference
timings on a current Fedora/x86\_64 machine with Python 3.15: generation about one second, full
verification about sixty-five seconds (dominated by the $15{,}960$-class dense recomputation),
independent checks under two seconds---so an outside reviewer can reproduce everything without
specialized hardware. A SHA-256 manifest covers the package;
it certifies byte integrity, not mathematical correctness.

\section{The failure at the preceding parameter}

At $\ell=5$ the window touches low degrees and the cancellation is no longer favourable: exact
computation gives central coefficients $(15,16,15,16,15)$ at degrees $18,\dots,22$, so the sequence
rises, falls, rises, falls, and unimodality fails while nonnegativity survives. This single instance
is the content of Theorem~\ref{thm:sharp}; it also explains the shape of the finite base: the margin
is smallest precisely at the bottom of the admissible range.

\section{Relationship to previous work}\label{sec:prior}

Unimodality of $\gauss{n}{k}$ goes back to Sylvester \cite{sylvester1878}; O'Hara \cite{ohara}
gave the first fully constructive proof, followed by Zeilberger's shorter exposition
\cite{zeilberger}; the KOH decomposition underlying both is the standard structural tool. Strict
unimodality results of Pak--Panova \cite{pakpanova} and Zanello \cite{zanello2015} give lower bounds
of magnitude $1$ on central first differences, insufficient for Theorem~\ref{thm:main}, where $3$ is
required at $j=17$ and $7$ at $j=11$. Reiner--Stanton \cite{reinerstanton} conjectured unimodality
of symmetric differences of Gaussian polynomials, of which the Stanley--Zanello family is a
generalization; their proved cases do not cover our slice. Koutschan, Uncu and Wong \cite{kuw}
resolved, using cylindrical algebraic decomposition and split points, two additional cases of the
Stanley--Zanello conjecture --- lower parameters $6$ and $7$, with explicit exception tables --- and
their $d$-strict monotonicity framework for single Gaussian polynomials stops computationally at
$m=7$; the $k=8$ slice treated here lies outside their proved coverage, as recorded with explicit
hypothesis/conclusion comparisons in the accompanying verification package. Quasipolynomiality of
restricted partition functions is classical (Sylvester; see \cite[Ch.~4]{stanleyec1}); in the
vector-partition setting it is a special case of the general theory of \cite{sturmfels}, and
positivity certificates based on finite exact computations are
standard in computer-assisted proofs; we claim no novelty for those tools as such.

\section{Limitations and open questions}\label{sec:open}

\begin{enumerate}
\item Scope. Theorem~\ref{thm:main} is a single fixed-parameter slice. It is not the full $k=8$
case of \cite{sz}, and it says nothing about odd $k$.
\item Certificate versus conceptual proof. Positivity of the margin polynomials is established by a
$15{,}960$-class exact certificate. A hand-checkable argument — for instance a closed form for the
residual quintic together with a sign analysis, or a combinatorial surplus injection of size $3$ at
$j=17$ — is missing and is the most natural next target.
\item Degree drop. The observed drop from degree $7$ to degree $5$ is uniform across all classes but
is presently recorded only as a computation; a structural explanation (cancellation between the
flat-top and right-edge layers of Section~\ref{sec:obstruction}) is open.
\item Sharpness of the tail. The tail margins are enormous while the finite base has margin $1$;
no optimality is claimed for $840$ as the junction between certificate and finite base. In fact the
truncation of Lemma~\ref{lem:ie} holds for every $\ell\ge1$ (the offending arguments are $\le -11$ uniformly),
while polynomiality of the surviving terms needs only $\ell\ge25$; a sharper packaging would certify
$\ell\ge25$ per residue class and shrink the finite verification to $6\le\ell\le24$. The uniform
parametrization used here trades that tightening for a single polynomial per residue class.
\item Verification trust. The certificate is machine-verified exact-integer computation in a
standard-library implementation; it has been reproduced by two independent implementations, but no
interactive proof assistant formalization exists.
\item No external peer review has been performed.
\end{enumerate}

\section*{Reproducibility statement}
All computations are exact integer arithmetic in Python~3 with the standard library only, deterministic
and randomized-sample-checked through a second code path; commands, reference outputs, software
versions, and file hashes are recorded in \texttt{verification/EXPECTED\_RESULTS.md}; the complete
runnable package is published at
\url{https://github.com/GeWoHu/gaussian-polynomial-unimodality} under the
\texttt{verification/} directory.

\bibliographystyle{plainnat}
\bibliography{references}

\end{document}
